\documentclass[EEEE]{KGCOEReport}

\usepackage[backend=biber]{biblatex}
\usepackage{graphicx}
\usepackage{newtxtext, newtxmath}
\usepackage{csquotes}
\usepackage{parskip}
\usepackage{karnaugh-map}
\usepackage{amsmath}
\usepackage{circuitikz}
\usepackage{hyperref}
\usepackage{float}
\usepackage{pdfpages}
\usepackage{pgfplots, pgfplotstable}

\graphicspath{./assets/}
\addbibresource{./assets/bib.bib}

\newcommand{\classCode}{EEEE 281}
\newcommand{\className}{Circuits 1}

\newcommand{\name}{Aden Perry}
% \newcommand{\partnerName}{PARTNER}

\newcommand{\LabSectionNum}{L7}

\newcommand{\TAs}{Sol Holston}
\newcommand{\exerciseNumber}{3}
\newcommand{\exerciseDescription}{Thevenin Equivalent Circuits and Linearity}

\newcommand{\dateDone}{2026-03-24}
\newcommand{\dateSubmitted}{2026-03-31}


\begin{document}

\maketitle

\section*{Abstract}


\section{Introduction and Theory}


\subsection{Theory: Circuit Investigated}


\subsection{Theory: LTspice Simulation}

\begin{table}[H]
	\center
	\caption{Theory}
	\begin{tabular}{|c|c|c|}
		\hline
		$R_L$ $\Omega$ & $V_{OUT}$ V & $I_{OUT}$ A \\
		\hline
		\hline
		2200           & 1.95556     & 0.000888889 \\
		\hline
		1000           & 1.90476     & 0.00190476  \\
		\hline
		470            & 1.80769     & 0.00384615  \\
		\hline
		100            & 1.33333     & 0.01333333  \\
		\hline
		68             & 1.15254     & 0.0169492   \\
		\hline
		47             & 0.969072    & 0.0206186   \\
		\hline
	\end{tabular}
\end{table}

\pgfplotstableread{
	X Y

	0.000888889	1.95556
	0.00190476 	1.90476
	0.00384615 	1.80769
	0.01333333 	1.33333
	0.0169492  	1.15254
	0.0206186  	0.969072
}\theorytbl

\begin{tikzpicture}
	\begin{axis}[legend pos=outer north east]
		\addplot [only marks, mark = *] table {\theorytbl};
		\addplot [thick, red] table[
				y={create col/linear regression={y=Y}}
			] % compute a linear regression from the input table
			{\theorytbl};
		\addlegendentry{$y(x)$}
		\addlegendentry{%
			$\pgfmathprintnumber{\pgfplotstableregressiona} \cdot x
				\pgfmathprintnumber[print sign]{\pgfplotstableregressionb}$}
	\end{axis}
\end{tikzpicture}

Theoretical $R_{TH} = 50$ ; Theoretical $V_{TH} = 2$



\section{Hardware Experiment: Results and Discussion}

\begin{table}[H]
	\center
	\caption{Hardware}
	\begin{tabular}{|c|c|c|}
		\hline
		$R_L$ $\Omega$ & $V_{OUT}$ V & $I_{OUT}$ A            \\
		\hline
		\hline
		2200           & 1.9590      & 0.00089045 \\
		\hline
		1000           & 1.9056      & 0.00190560 \\
		\hline
		470            & 1.8053      & 0.00384106 \\
		\hline
		100            & 1.3166      & 0.01316600 \\
		\hline
		68             & 1.1353      & 0.01669559 \\
		\hline
		47             & 0.95256     & 0.02026723 \\
		\hline
	\end{tabular}
\end{table}

\pgfplotstableread{
	X Y

	0.00089045454545454551 1.9590
	0.0019055999999999999 1.9056
	0.0038410638297872337 1.8053
	0.013166000000000001  1.3166
	0.016695588235294118 1.1353
	0.020267234042553192 0.95256

}\hwtbl

\begin{tikzpicture}
	\begin{axis}[legend pos=outer north east]
		\addplot [only marks, mark = *] table {\hwtbl};
		\addplot [thick, red] table[
				y={create col/linear regression={y=Y}}
			] % compute a linear regression from the input table
			{\hwtbl};
		\addlegendentry{$y(x)$}
		\addlegendentry{%
			$\pgfmathprintnumber{\pgfplotstableregressiona} \cdot x
				\pgfmathprintnumber[print sign]{\pgfplotstableregressionb}$}
	\end{axis}
\end{tikzpicture}

Hardware $R_{TH} = 52.03$ ; Hardware $V_{TH} = 2$

\section{Conclusion}

Having a positive (non-zero) Thevenin resistance decreases the output voltage
from the input voltage.

\section{Acknowledgements}

The template for this Technical Memo was modified from a \LaTeX \ class created
by Joshua Milas and Seth Gower, using the provided Tech Memo Template and
example Tech Memo from the Advanced Course as inspiration.

The Teaching Assistant, Sol Holston, verified that calculated, simulated, and
measured values were correct.

\printbibliography

\end{document}
